problem:  
Let \( (M^4, g_0) \) be a smooth, closed, oriented 4‑manifold with positive isotropic curvature.  
Assume the existence of a sequence of **Ricci flows with surgery** \( (M, g_k(t))_{t\ge0} \) starting at \(g_0\), with surgery parameters \(\delta_k \to 0\), such that the flows are non‑collapsing and converge (after passing to a subsequence) in the space–time Cheeger–Gromov sense to a singular Ricci flow \( (\mathcal{X}, g_\infty(t)) \) on \(t>0\).  

Let \((\mathcal{X}_{\mathrm{reg}}, g_\infty(t))\) denote the regular part, and \(\mathcal{S} = \mathcal{X} \setminus \mathcal{X}_{\mathrm{reg}}\) the singular set. For each \(t>0\), define the rescaled metrics  
\[
h_\lambda(t) = \lambda\, g_\infty(\lambda^{-1}t),
\]
and suppose that for a sequence \(\lambda_j \to \infty\) the pointed flows \((\mathcal{X}, h_{\lambda_j}(t), x_j)\), based at points \(x_j \in \mathcal{S}\), subconverge to complete ancient solutions \( (\mathcal{Y}, h_\infty(t)) \) with nonnegative curvature operator.  

**Question.** —  
Must any such tangent flow \( (\mathcal{Y}, h_\infty(t)) \) arising from blow‑up around a singular point \(x_j \to x_\infty \in \mathcal{S}\) be a **gradient shrinking Ricci soliton**?  
Equivalently, does every singularity in the vanishing‑surgery limit of 4‑dimensional Ricci flow possess a canonical solitonic tangent modeled on a self‑similar shrinker, or can one construct (within this limiting framework) a non‑solitonic, nonselfsimilar ancient tangent flow indicating a fundamentally new singularity type in dimension four?

Why it is a "good" problem:  
This problem targets the uncharted structure of **singularities in 4‑dimensional weak Ricci flows**, asking whether all tangent flows are necessarily gradient shrinking solitons—the higher‑dimensional analog of Perelman’s canonical neighborhood conjecture in 3D.  

Its statement is compact and geometrically natural: “Are 4D singular tangent flows always shrinkers?” Yet answering it would require developing or refuting the extension of one of Ricci flow’s deepest simplifications—singular self‑similarity—beyond the presently settled three‑dimensional case.  

A positive resolution would solidify the universality of self‑similarity as the canonical asymptotic structure of singularities, thus providing analytic control and potential classification of 4D singularities. A negative resolution would reveal the existence of truly **non‑solitonic ancient solutions** in 4D, signaling that the analytic universe of Ricci flow singularities is vastly richer than expected.  

The problem connects geometry (structure of singularities), analysis (classification of ancient solutions), and topology (global implications for 4‑manifolds). It distills open challenges in the long‑time and singularity behavior of Ricci flow into a single clear conjectural dichotomy—hence a research‑level “narrow entrance, profound depth” question.